\subsection{对固有态射的应用：I．交换性质}
\source{第三卷第二部分，原印刷第65--72页}
接下来的三个目，主要是把前面各目的结果译写为预概形态射的语言。

\egap{7.7.1}
设 $f:X\to Y$ 为拟紧且分离的预概形态射，$\PP_\bullet$ 为拟凝聚 $\OO_X$-模复形，微分次数为 $-1$，并设各 $\PP_i$ 在 $Y$ 上平坦（$0_{\mathrm I}$，6.7.1）。考虑拟凝聚 $\OO_Y$-模范畴中的 $\delta$-函子 $\M\mapsto\TT_\bullet(\M)$，也记作 $\TT_\bullet(\PP_\bullet,\M)$，定义为
\begin{equation}
 \TT_n(\PP_\bullet,\M)=\TT_n(\M)
 =\cH^{-n}(f,\PP^\bullet\otimes_{\OO_Y}\M),\qquad n\in\mathbf Z.
 \tag{7.7.1.1}
\end{equation}
这里 $\PP^j=\PP_{-j}$，微分次数因而为 $+1$。由（6.2.3），其值为拟凝聚 $\OO_Y$-模；由（6.2.6），它是关于 $\M$ 的同调函子。

\egap{7.7.2}
设 $g:Y'\to Y$ 为态射，置 $X'=X_{(Y')}=X\times_Y Y'$，$f'=f_{(Y')}:X'\to Y'$，它仍拟紧且分离。置 $\PP'_\bullet=\PP_\bullet\otimes_{\OO_Y}\OO_{Y'}$；由（I，9.1.12）和（$0_{\mathrm I}$，6.2.1），其各项拟凝聚且在 $Y'$ 上平坦。采用同样的次数约定，置
\begin{equation}
 \begin{aligned}
 \TT'_\bullet(\M')
 &=\TT_\bullet(f',\PP'_\bullet\otimes_{\OO_{Y'}}\M')\\
 &=\cH^{-\bullet}(f',\PP'^\bullet\otimes_{\OO_{Y'}}\M').
 \end{aligned}\tag{7.7.2.1}
\end{equation}
这是关于拟凝聚 $\OO_{Y'}$-模 $\M'$ 的同调函子。当 $Y'=\Spec A'$ 时，以 $\TT^{A'}$ 代替 $\TT'$；对每个 $A'$-模 $M'$，置
\[
 T^{A'}(M')=\Gamma(Y',\TT^{A'}(\widetilde{M'})),
 \qquad\TT^{A'}(\widetilde{M'})=\widetilde{T^{A'}(M')}.
\]
$T^{A'}$ 为取值于 $A'$-模范畴的同调函子。

若 $Y=\Spec A$ 也仿射，则 $T^{A'}$ 与从 $T^A$ 标量扩张到 $A'$ 所得函子相同（7.1.3）。事实上，若 $g':X'\to X$ 为相应投影，则它仿射（II，1.6.2）；仿射开覆盖 $\mathfrak U$ 的逆像 $\mathfrak U'=g'^{-1}(\mathfrak U)$ 仍为仿射开覆盖。由（6.2.2），只须验证
\[
 C^\bullet(\mathfrak U,\PP_\bullet\otimes_{\OO_Y}g_*(\M'))
 =C^\bullet(\mathfrak U',\PP'_\bullet\otimes_{\OO_{Y'}}\M').
\]
最终归结为对仿射开集 $U\subseteq X$ 及 $U'=g'^{-1}(U)$，
\[
 \Gamma(U,\PP_\bullet\otimes_{\OO_Y}g_*(\M'))
 =\Gamma(U',\PP'_\bullet\otimes_{\OO_{Y'}}\M'),
\]
这由（I，1.3、3.2）直接得到。特别地，对 $Y$ 的开集 $U$ 及任意拟凝聚 $\M$，
\begin{equation}
 \TT^U_\bullet(\M|U)=(\TT_\bullet(\M))|U.\tag{7.7.2.2}
\end{equation}

\egap{7.7.3}
对每个拟凝聚 $\OO_Y$-模 $\M$，有关于 $\M$ 函子性的典范同态
\begin{equation}
 \TT_p(\OO_Y)\otimes_{\OO_Y}\M\longrightarrow\TT_p(\M).
 \tag{7.7.3.1}
\end{equation}
当 $Y$ 仿射时，这已经在（7.2.2）定义。一般情形通过限制相容性得到：对仿射开集 $V\subseteq U\subseteq Y$，令
\[
 A_W=\TT_p^W(\OO_W)\otimes_{\OO_W}(\M|W),\qquad
 B_W=\TT_p^W(\M|W)\quad(W=U,V),
\]
则在与限制层作典范识别后，图
\[
 \begin{tikzcd}
 A_U\arrow[r]\arrow[d,"\text{限制}"']&B_U\arrow[d,"\text{限制}"]\\
 A_V\arrow[r]&B_V
 \end{tikzcd}
\]
由（7.2.3.3）交换。

对任一态射 $g:Y'\to Y$，还有典范同态
\begin{equation}
 \TT_p(\OO_Y)\otimes_{\OO_Y}\OO_{Y'}\longrightarrow
 \TT'_p(\OO_{Y'}).\tag{7.7.3.2}
\end{equation}
它是（6.7.11.2）关于终项的特殊情形：取 $S=Y$、$u_1=f$、$v_1=1_Y$，并把第二个复形置为次数0的模。若 $Y=\Spec A$、$Y'=\Spec A'$ 均仿射，则（7.7.3.2）与（7.2.2）中典范 $A'$-线性同态
\[
 T_p^A(A)\otimes_A A'\longrightarrow T_p^{A'}(A')=T_p^A(A')
\]
相伴。这可从（6.7.11）看出，因为可在那里选取 $\mathfrak U'^{(i)}=u_i^{-1}(\mathfrak U^{(i)})$。

\egap{7.7.4}
当 $f$ 固有、$Y=\Spec A$ 仿射诺特、$\PP_\bullet$ 下有界且由凝聚、$Y$-平坦的模组成时，由（6.10.5），在同构意义下可写
\[
 \TT_\bullet(\M)=\HH_\bullet(\LL_\bullet\otimes_{\OO_Y}\M),
 \qquad\LL_\bullet=\widetilde{L_\bullet},
\]
其中 $L_\bullet$ 为下有界的有限型自由 $A$-模复形。因此 $T^A$ 正属于（7.4）、（7.6）详细研究的类型。现在译写该研究的结果。

\egas{7.7.5}{定理}
设 $Y$ 为局部诺特预概形，$(U_\alpha)$ 为其仿射开覆盖，$f:X\to Y$ 固有，$\PP_\bullet$ 为下有界的凝聚、$Y$-平坦 $\OO_X$-模复形。由（7.7.1.1）定义的同调函子具有以下性质。

\textbf{I．半连续性。}函数
\begin{equation}
 y\longmapsto d_p(y)=\dim_{k(y)}T_p^{k(y)}(k(y))\tag{7.7.5.1}
\end{equation}
上半连续。
\orig{本定理的一个特殊情形已经出现在 Chow--Igusa 的札记 [3] 中。对解析空间，在较特殊的假设下，Kodaira--Spencer 已证明半连续性（\textit{On the variations of almost-complex structures}，载 \textit{Algebraic Geometry and Topology, A Symposium in honor of S. Lefschetz}，Princeton 数学丛书第12册，第139--150页，Princeton，1957年）。一般版本由 Grauert [31] 证明。}

\textbf{II．交换性质。}给定整数 $p$，下列条件等价：
\begin{enumerate}[label=\alph*)]
\item $\TT_p$ 右正合。
\item[\textup{a')}] $\TT_p(\M)$ 同构于某个形如 $\N\otimes_{\OO_Y}\M$ 的函子；此时 $\N$ 必同构于 $\TT_p(\OO_Y)=\cH^{-p}(f,\PP^\bullet)$。
\item[\textup{a'')}] 典范函子同态（7.7.3.1）为同构。
\item[\textup{b)}] $\TT_{p-1}$ 左正合。
\item[\textup{b')}] 存在 $\OO_Y$-模 $\mathscr Q$（必凝聚，并确定至唯一同构）及函子同构
\begin{equation}
 \TT_{p-1}(\M)\simeq\shHom_{\OO_Y}(\mathscr Q,\M).
 \tag{7.7.5.2}
\end{equation}
\item[\textup{c)}] 若 $A_\alpha$ 为 $U_\alpha$ 的环，则每个 $T_p^{A_\alpha}$ 都右正合。
\item[\textup{d)}] 对每个态射 $g:Y'\to Y$，典范同态
\begin{equation}
 \TT_p(\OO_Y)\otimes_{\OO_Y}\OO_{Y'}\longrightarrow\TT'_p(\OO_{Y'})
 \tag{7.7.5.3}
\end{equation}
为同构。
\end{enumerate}
半连续性是 $Y$ 上的局部性质，故来自（7.7.4）和（7.6.9）。a'') 蕴含 a')，a') 蕴含 a)。$Y$ 仿射时，由（7.7.4），a)、a'')、b)、b') 的等价性已经由（7.3.1）、（7.4.6）证明。

为推广到一般情形，先证明 a) 与 c) 等价，这也说明 a) 在 $Y$ 上局部。相同论证适用于 a'')、b)。c) 蕴含 a)；反之，只须对任一仿射开集 $U\subseteq Y$ 及拟凝聚 $\OO_U$-模短正合列
\[
 0\longrightarrow\F'\longrightarrow\F\longrightarrow\F''\longrightarrow0
\]
找出 $Y$ 上的拟凝聚模短正合列，其限制为该列。由 $Y$ 局部诺特及（I，9.4.2），先把 $\F$ 延拓为 $\G$，把 $\F'$ 延拓为 $\G$ 的拟凝聚子模 $\G'$，再取 $\G''=\G/\G'$ 即可。

对一般 $Y$，证明 b)、b') 等价。仿射情形中的表示模 $\mathscr Q$ 确定至唯一同构；因此对仿射 $Y$ 的仿射开集 $U$，其表示模就是 $\mathscr Q|U$。一般地，在每个仿射开集 $U$ 上取凝聚表示模 $\mathscr Q_U$。若仿射开集 $V\subseteq U$，则 $\mathscr Q_U|V$ 与 $\mathscr Q_V$ 典范同构；故可粘合，得到满足（7.7.5.2）的唯一表示模。

最后证明 a)、d) 等价。二者在 $Y$ 上均局部，且 d) 在 $Y'$ 上也局部。故设 $Y=\Spec A$、$Y'=\Spec A'$。已知 $T^{A'}$ 由 $T^A$ 标量扩张得到，a') 便给出（7.7.5.3）为同构。反之，设 d) 成立。对任一 $A$-模 $M$，取平方零扩张 $A'=A\oplus M$，乘法为
\[
 (a_1,m_1)(a_2,m_2)=(a_1a_2,a_1m_2+a_2m_1).
\]
此时 $T_p(A')=T_p(A)\oplus T_p(M)$；（7.7.5.3）为双射，推出典范同态 $T_p(A)\otimes_A M\to T_p(M)$ 也为双射，即 a'')。

\egas{7.7.6}{定理}
设 $Y$ 局部诺特，$f:X\to Y$ 固有，$\F$ 为凝聚且 $Y$-平坦的 $\OO_X$-模。则存在凝聚 $\OO_Y$-模 $\mathscr Q$（确定至唯一同构），使对拟凝聚 $\OO_Y$-模 $\M$ 有函子同构
\begin{align}
 f_*(\F\otimes_{\OO_Y}\M)&\simeq\shHom_{\OO_Y}(\mathscr Q,\M),\tag{7.7.6.1}\\
 \Gamma(X,\F\otimes_{\OO_Y}\M)&\simeq\Hom_{\OO_Y}(\mathscr Q,\M).
 \tag{7.7.6.2}
\end{align}
因 $\M\mapsto\F\otimes_{\OO_Y}\M$ 正合（$0_{\mathrm I}$，6.7.4），$f_*$ 左正合，复合函子左正合。取 $p=1$，应用（7.7.5，b）与 b') 的等价性即可。

\egas{7.7.7}{推论}
设 $Y$ 局部诺特，$f:X\to Y$ 固有，$\F,\F'$ 凝聚且 $Y$-平坦，$u:\F\to\F'$ 为同态。对拟凝聚 $\OO_Y$-模 $\M$，置
\[
 \begin{aligned}
 \TT(\M)&=\Ker\bigl(f_*(\F\otimes\M)\to f_*(\F'\otimes\M)\bigr),\\
 T(\M)&=\Gamma(Y,\TT(\M))\\
 &=\Ker\bigl(\Gamma(X,\F\otimes\M)\to\Gamma(X,\F'\otimes\M)\bigr).
 \end{aligned}
\]
则存在凝聚 $\OO_Y$-模 $\B$（确定至唯一同构），使
\begin{align}
 \TT(\M)&\simeq\shHom_{\OO_Y}(\B,\M),\tag{7.7.7.1}\\
 T(\M)&\simeq\Hom_{\OO_Y}(\B,\M).\tag{7.7.7.2}
\end{align}
只须证明（7.7.7.2）：它在仿射情形给出（7.7.7.1），再像（7.7.5）的 b)、b') 那样，利用表示对象的唯一性（0，8.1.8）粘合。由（7.7.6），存在凝聚 $\mathscr Q,\mathscr Q'$，分别表示 $\Gamma(X,\F\otimes-)$、$\Gamma(X,\F'\otimes-)$。$u$ 定义这两个函子之间的态射，由（0，8.1.4），唯一对应于 $v:\mathscr Q'\to\mathscr Q$，使图
\[
 \begin{tikzcd}
 \Gamma(X,\F\otimes\M)\arrow[r]\arrow[d,"\sim"']&\Gamma(X,\F'\otimes\M)\arrow[d,"\sim"]\\
 \Hom_{\OO_Y}(\mathscr Q,\M)\arrow[r]&\Hom_{\OO_Y}(\mathscr Q',\M)
 \end{tikzcd}
\]
交换。反变函子 $\N\mapsto\Hom_{\OO_Y}(\N,\M)$ 左正合，故取 $\B=\Coker v$ 即可。

\egas{7.7.8}{推论}
设 $X,Y,f$ 满足（7.7.6）的假设，$\F,\G$ 为凝聚 $\OO_X$-模，满足：（i）$\F$ 在 $Y$ 上平坦；（ii）$\G$ 同构于有限型局部自由模之间某同态 $v:\E_1\to\E_0$ 的余核。置
\[
 \TT(\M)=f_*\shHom_{\OO_X}(\G,\F\otimes_{\OO_Y}\M),\qquad
 T(\M)=\Hom_{\OO_X}(\G,\F\otimes_{\OO_Y}\M).
\]
则存在凝聚 $\OO_Y$-模 $\N$（确定至唯一同构），使
\begin{align}
 \TT(\M)&\simeq\shHom_{\OO_Y}(\N,\M),\tag{7.7.8.1}\\
 T(\M)&\simeq\Hom_{\OO_Y}(\N,\M).\tag{7.7.8.2}
\end{align}
由（$0_{\mathrm I}$，5.4.2.1），对 $i=0,1$ 有函子性同构
\[
 \begin{aligned}
 \shHom_{\OO_X}(\E_i,\F\otimes\M)
 &\simeq\E_i^\vee\otimes_{\OO_X}(\F\otimes\M)\\
 &\simeq(\E_i^\vee\otimes_{\OO_X}\F)\otimes\M
 \simeq\shHom_{\OO_X}(\E_i,\F)\otimes\M.
 \end{aligned}
\]
令 $\F_i=\shHom_{\OO_X}(\E_i,\F)$；它们凝聚（$0_{\mathrm I}$，5.3.5）且 $Y$-平坦（同章5.4.2）。置 $u=\shHom(v,1_\F):\F_0\to\F_1$。由 Hom 函子的左正合性，
\[
 \shHom_{\OO_X}(\G,\F\otimes\M)
 \simeq\Ker(\F_0\otimes\M\to\F_1\otimes\M).
\]
再因 $f_*$ 左正合，$\TT(\M)$ 同构于 $\Ker(f_*(\F_0\otimes\M)\to f_*(\F_1\otimes\M))$，应用（7.7.7）即可。

\egas{7.7.9}{注记}
（i）在（7.7.6）—（7.7.8）中，$\mathscr Q,\B,\N$ 的构造与基变换交换。例如采用（7.7.2）的记号，在（7.7.6）情形，对拟凝聚 $\OO_{Y'}$-模 $\M'$ 有
\[
 f'_*(\F'\otimes_{\OO_{Y'}}\M')
 \simeq\shHom_{\OO_{Y'}}(g^*\mathscr Q,\M').
\]
由（7.7.2）的说明，归结为伴随同构
\[
 \Hom_{\OO_Y}(\mathscr Q,g_*\M')
 =\Hom_{\OO_{Y'}}(g^*\mathscr Q,\M')
\]
（$0_{\mathrm I}$，4.4.3.1）。同样，在（7.7.7）的基变换中用 $g^*\B$ 代替 $\B$；在（7.7.8）中用 $g^*\N$ 代替 $\N$，因为基变换后的 $\E'_1\to\E'_0\to\G'\to0$ 仍正合，并且
\[
 \shHom_{\OO_{X'}}(\E'_i,\F\otimes\OO_{Y'})
 \simeq\E_i^\vee\otimes_{\OO_X}\F\otimes_{\OO_Y}\OO_{Y'}.
\]

（ii）若存在 $Y$-丰沛可逆 $\OO_X$-模，例如 $Y$ 仿射且 $f$ 射影，则任意凝聚 $\G$ 均满足（7.7.8，ii）。结合（II，5.5.1），由（II，2.7.10）可取有限型局部自由 $\E_0$ 满射到 $\G$；其核凝聚，对核再次应用同一结果，即得所需表示。

（iii）我们将在第V章证明，（7.7.8）的限制性假设（ii）是多余的。

\egas{7.7.10}{命题（交换性质的局部判据）}
在（7.7.5）的一般条件下，给定 $y\in Y$ 和整数 $p$，以下性质等价：
\begin{enumerate}[label=\alph*)]
\item $T_p^{\OO_y}$ 右正合。
\item 典范同态 $T_p^{\OO_y}(\OO_y)\to T_p^{\OO_y}(k(y))$ 满射。
\item 对每个整数 $n$，典范同态 $T_p^{\OO_y}(\OO_y/\mathfrak m_y^{n+1})\to T_p^{\OO_y}(k(y))$ 满射。
\end{enumerate}
此外，满足这些条件的点构成使 $\TT_p^U$ 右正合的最大开集 $U\subseteq Y$。

结合（7.7.4），等价性来自（7.4.7）及（7.5.2）；开性来自（7.4.4）。反之，若 $\TT_p^U$ 右正合，则由（7.7.5，c）及（7.6.1），$y\in U$ 时 $T_p^{\OO_y}$ 右正合。

\egas{7.7.11}{推论}
若 $\TT_p$ 右正合（分别左正合），则任意基变换 $g:Y'\to Y$ 后的 $\TT'_p$ 也如此；若 $g$ 忠实平坦，则逆命题成立。

第一断言来自（7.6.1）以及正合性对 $Y,Y'$ 局部（7.7.5，c、b）。反之，由（7.7.10），只须在每个 $y\in Y$ 证明 $T_p^{\OO_y}$ 的相应正合性。因 $g$ 忠实平坦，存在 $y'$ 映到 $y$，且 $\OO_{y'}$ 在 $\OO_y$ 上忠实平坦；结论来自（7.6.1）。

\egas{7.7.12}{注记}
（i）在（7.7.4）的假设下，再设 $\PP_\bullet$ 为有限复形。取 $X$ 的有限仿射开覆盖 $\mathfrak U$，则 $C^\bullet(\mathfrak U,\PP^\bullet\otimes\M)$ 为有限双复形。准确地说，存在与 $\M$ 无关的有限指标集 $E$，使所有 $(h,k)\notin E$ 的项都为零。因此存在 $i_0$，使 $i\geq i_0$ 时对任意拟凝聚 $\M$ 都有 $\TT_i(\M)=0$。这些函子正合，故相应的 $Z'_{i_0}(L_\bullet)$ 为有限型平坦、因而有限型投射 $A$-模。取
\[
 L'_i=L_i\ (i<i_0),\qquad L'_{i_0}=Z'_{i_0}(L_\bullet),
 \qquad L'_i=0\ (i>i_0),
\]
并令 $\LL'_\bullet=\widetilde{L'_\bullet}$。显然 $i<i_0-1$ 时两复形张量 $\M$ 后的同调相同，$i\geq i_0$ 时也相同（均为零）；$i=i_0-1$ 时，由 $\Im(Z'_{i_0}\otimes_A M\to L_{i_0-1}\otimes_A M)=\Im(L_{i_0}\otimes_A M\to L_{i_0-1}\otimes_A M)$，同调亦相同。故在（7.7.4）中，若只要求 $\LL_i$ 有限型局部自由，即对应有限型投射 $A$-模，就可令 $\LL_\bullet$ 为有限复形。

特别地，若 $\PP_\bullet$ 只在次数0有一非零项 $\F$，则 $\TT_p(\M)=R^{-p}f_*(\F\otimes\M)$，可令 $i>0$ 时 $\LL_i=0$；这时宜采用上同调记号，以 $\TT^{-p}$ 代替 $\TT_p$。

（ii）在（7.7.5）中若去掉各 $\PP_i$ 的 $Y$-平坦性，结论在改定义为
\begin{equation}
 \TT_p(\M)=\lTor_p^Y(f,1_Y;\PP_\bullet,\M)\tag{7.7.12.1}
\end{equation}
后仍成立。由（6.7.9），$\lTor_p^Y(f,1_Y;\PP_\bullet,\OO_Y)$ 凝聚；结合（6.10.1），（6.10.5）的证明不变，仍在仿射底空间上得到 $\TT_\bullet(\M)=\HH_\bullet(\LL_\bullet\otimes\M)$，其中 $\LL_\bullet$ 与有限型自由模复形相伴。这就证明了断言。
